\documentclass[12pt,a4paper]{article}
\usepackage{fontspec}\usepackage{polyglossia}\setdefaultlanguage{ukrainian}
\setmainfont{NewComputerModern10}\setsansfont{Noto Sans}\setmonofont{Noto Sans Mono}[Scale=.82]
\usepackage{amsmath,amssymb,booktabs,caption,enumitem,float,geometry,graphicx,hyperref,microtype,titlesec,xcolor}
\geometry{left=2cm,right=2cm,top=2cm,bottom=2cm}\setlength{\parindent}{1.25cm}\linespread{1.08}
\titleformat{\section}{\normalfont\fontsize{14}{16}\selectfont\bfseries}{\thesection}{.6em}{}\titlespacing*{\section}{0pt}{1.3em}{.6em}
\captionsetup{font=small,labelfont=bf,labelsep=endash}
\usepackage[newfloat]{minted}\SetupFloatingEnvironment{listing}{name=Лістинг}
\setminted{fontsize=\footnotesize,linenos,breaklines,autogobble,frame=single,framesep=2mm,bgcolor=gray!5,tabsize=4}
\hypersetup{unicode,colorlinks=true,linkcolor=black,pdftitle={ІО-41, Давидчук А. М., Лабораторна робота №8},pdfauthor={Давидчук А. М.}}
\newcommand{\resultDir}{../results}\newcommand{\studentGroup}{ІО-41}\newcommand{\studentName}{Давидчук А. М.}\newcommand{\reportYear}{2026}
\begin{document}\pagenumbering{gobble}
\begin{titlepage}
\begin{center}\includegraphics[width=.7\textwidth]{../../../../../../templates/letterhead.png}\end{center}
\begin{center}\large Національний технічний університет України\\ «Київський політехнічний інститут імені Ігоря Сікорського»\\[1em]Факультет інформатики та обчислювальної техніки\\Кафедра обчислювальної техніки\end{center}
\vfill\begin{center}\textbf{\LARGE Методи та технології штучного інтелекту}\\[2em]\textbf{\Large Лабораторна робота №8}\\[.7em]«Застосування генетичних алгоритмів у задачах оптимізації»\end{center}
\vfill\begin{flushright}\begin{minipage}{.48\textwidth}Виконав: студент 3 курсу ФІОТ,\\групи \studentGroup\\\textit{\studentName}\\[1em]Перевірив: \textit{Польшакова О. М.}\end{minipage}\end{flushright}
\vfill\begin{center}Київ --- \reportYear\end{center}
\end{titlepage}\pagenumbering{arabic}

\section{Тема та мета роботи}
\noindent\textbf{Тема:} генетична оптимізація структури обчислювальної мережі.\\
\noindent\textbf{Мета:} сформувати модель мережі, закодувати підключення вузлів хромосомою, розробити цільову функцію й генетичний алгоритм, який зменшує нормоване навантаження на магістраль за обмежень місткості комутаторів і довжини абонентських ліній.

\section{Вибір індивідуального завдання}
Номер студента 6 є парним, тому відповідно до методички виконується задача оптимізації обчислювальної мережі. Розглядаються 24 робочі станції та чотири послідовно з'єднані комутатори. До кожного комутатора можна підключити не більше шести станцій.

Для кожної пари станцій задано середній обсяг обміну даними. Частина станцій утворює логічні робочі групи з інтенсивним внутрішнім трафіком. Якщо станції під'єднані до різних комутаторів, їхній обмін проходить магістральними каналами. Завдання полягає у такому перепідключенні, щоб активні пари були ближче одна до одної, але довжина кабелів до комутаторів не стала надмірною.

\section{Математична модель}
Хромосома є цілочисловим вектором
\begin{equation}C=[c_1,c_2,\ldots,c_{24}],\qquad c_i\in\{0,1,2,3\},\end{equation}
де $c_i$ --- номер комутатора, до якого підключено вузол $i$. Обмеження місткості вимагає, щоб кожне значення зустрічалося рівно шість разів.

Для пари вузлів $i,j$ з інтенсивністю трафіку $B_{ij}$ кількість пройдених магістральних сегментів дорівнює $|c_i-c_j|$. Сумарна складова магістрального трафіку:
\begin{equation}T(C)=\sum_{i<j}B_{ij}|c_i-c_j|.\end{equation}
Довжина абонентських ліній дорівнює
\begin{equation}L(C)=\sum_i d(u_i,q_{c_i}),\end{equation}
де $d$ --- евклідова відстань між робочою станцією і комутатором. Відповідно до методички ефективність визначається відношенням трафіку каналу $T_k$ до його пропускної здатності $P_k$. Для трьох сегментів магістралі прийнято $P_k=250$ умовних одиниць. Додатково задано максимально допустиму довжину одного абонентського кабелю $l_{\max}=55$. Мінімізується функція
\begin{equation}
F(C)=\sum_{k=1}^{3}\frac{T_k(C)}{P_k}+10\sum_{i=1}^{24}\max\!\left(0,l_i(C)-l_{\max}\right)^2.
\end{equation}
Перший доданок зменшує завантаженість магістралі, а штраф робить схеми з надто довгими підключеннями невигідними. Якщо всі обмеження виконані, штраф дорівнює нулю.

\section{Параметри генетичного алгоритму}
Початкова популяція містить 80 випадкових допустимих схем. Велика порівняно з лабораторною №7 популяція потрібна через значно більший дискретний простір: без урахування симетрії кількість призначень дорівнює $4^{24}$. Еволюція триває 250 поколінь.

Для вибору батьків застосовано описаний у методичці пропорційний відбір. Оскільки задача є задачею мінімізації, оцінки масштабуються за формулою $w_i=F_{\max}-F_i$, після чого нормуються до ймовірностей $p_i=w_i/\sum_jw_j$. Отже, менше значення цільової функції дає більшу ймовірність розмноження. Імовірність одноточкового кросоверу становить 0,90. Після нього спеціальна процедура відновлює обмеження місткості. Мутація з імовірністю 0,15 міняє місцями комутатори двох вузлів, тому місткість автоматично зберігається. Чотири найкращі особини переносяться без змін. Критерій зупинки --- 250 поколінь; фіксований seed робить експеримент повторюваним.

\section{Програмна реалізація}
Цільова функція одночасно обчислює сумарний магістральний трафік, навантаження кожного каналу і довжину кабелів.
\begin{listing}[H]
\inputminted[firstline=42,lastline=70]{python}{../src/main.py}
\caption{Відновлення допустимої хромосоми та оцінювання підключення}
\end{listing}

Основний цикл ГА виконує ранжування, масштабування пристосованості, пропорційний відбір, елітизм, кросовер, repair і мутацію перестановкою.
\begin{listing}[H]
\inputminted[firstline=79,lastline=111]{python}{../src/main.py}
\caption{Генетична оптимізація мережі}
\end{listing}

\section{Результати}
За початкову схему взято підключення до найближчого комутатора з подальшим відновленням місткості. Результати порівняння наведено в таблиці.
\begin{table}[H]\centering
\begin{tabular}{lrr}\toprule Показник & До оптимізації & Після оптимізації\\\midrule
Цільова функція $F$ & 2,3860 & 0,8181\\
Сумарне нормоване навантаження & 2,3860 & 0,8181\\
Магістральна складова $T$ & 596,4997 & 204,5364\\
Сумарна довжина $L$ & 621,3469 & 716,7299\\
Максимальна довжина одного кабелю & 50,9382 & 51,7089\\
Порушення обмеження $l_i\leq55$ & 0 & 0\\
Покращення $F$ & \multicolumn{2}{c}{65,71\%}\\\bottomrule
\end{tabular}\caption{Результат оптимізації підключення}\end{table}

Генетичний алгоритм зменшив сумарне нормоване навантаження у 2,92 раза. Сумарна довжина кабелів збільшилася, однак найбільша окрема довжина становить 51,7089 і не перевищує допустимі 55 одиниць. Тому штрафовий доданок дорівнює нулю як до, так і після оптимізації.

\begin{figure}[H]\centering\includegraphics[width=.88\textwidth]{\resultDir/topology_before.png}\caption{Початкове підключення станцій}\end{figure}
Початкова схема переважно мінімізує фізичну відстань, але не враховує матрицю трафіку. Тому логічно пов'язані вузли можуть опинитися на протилежних ділянках магістралі.

\begin{figure}[H]\centering\includegraphics[width=.88\textwidth]{\resultDir/topology_after.png}\caption{Підключення після генетичної оптимізації}\end{figure}
Після оптимізації просторово віддалені станції іноді підключаються до одного комутатора. Це виправдано, якщо між ними відбувається інтенсивний обмін, а довжина кожного такого підключення залишається в допустимих межах.

\begin{figure}[H]\centering\includegraphics[width=.82\textwidth]{\resultDir/convergence.png}\caption{Збіжність генетичного алгоритму}\end{figure}
Найкраще значення швидко зменшується на початку, а далі покращується невеликими кроками. Середнє значення популяції також спадає, отже відбір поширює вдалі фрагменти схем, а мутація продовжує пошук нових комбінацій. Логарифмічна шкала дає змогу одночасно побачити великі штрафи недопустимих початкових схем і мале значення найкращої допустимої схеми.

\begin{figure}[H]\centering\includegraphics[width=.75\textwidth]{\resultDir/backbone_loads.png}\caption{Навантаження магістральних каналів до та після оптимізації}\end{figure}
Зменшення стовпців для всіх трьох сегментів підтверджує, що покращення цільової функції не зводиться до перенесення перевантаження на інший канал. Пунктирна лінія показує прийняту пропускну здатність 250 одиниць.

\section{Висновки}
Розв'язано комбінаторну задачу перепідключення 24 станцій до чотирьох комутаторів з обмеженням шість станцій на комутатор. Реалізовано пропорційний відбір, елітизм, одноточковий кросовер, відновлення допустимості та мутацію перестановкою. Генетичний алгоритм зменшив цільову функцію з 2,3860 до 0,8181, тобто на 65,71\%, а магістральну складову --- з 596,4997 до 204,5364. Обмеження максимальної довжини кабелю не порушено, тому отримана схема є допустимою.

\section{Контрольні питання}
\begin{enumerate}[leftmargin=*]
\item \textit{У чому сутність задач оптимізації, які розв'язують генетичними алгоритмами?}\\
Потрібно знайти значення змінних, що мінімізують або максимізують цільову функцію за заданих обмежень.
\item \textit{Чому ГА ефективні за великої кількості параметрів?}\\
Вони досліджують простір популяцією рішень і комбінують вдалі частини без повного перебору всіх варіантів.
\item \textit{Коли традиційні методи поступаються генетичним?}\\
За дискретних, розривних, нелінійних і багатомодальних функцій, для яких немає зручних похідних.
\item \textit{Як визначається функція пристосованості?}\\
Вона має чисельно відображати мету задачі та штрафувати порушення обмежень; у роботі це сума відношень $T_k/P_k$ плюс штраф за перевищення допустимої довжини кабелю.
\item \textit{Чому правильна fitness-функція критично важлива?}\\
Алгоритм оптимізує саме її, тому невдале формулювання може привести до математично кращого, але практично неприйнятного рішення.
\item \textit{Як ГА уникає локальних екстремумів?}\\
Популяція одночасно охоплює кілька областей, а мутація створює нові варіанти поза поточними комбінаціями.
\item \textit{Для яких типів оптимізації застосовують ГА?}\\
Для неперервних, дискретних, комбінаторних, нелінійних і багатокритеріальних задач.
\item \textit{Як кросовер допомагає пошуку?}\\
Він поєднує частини двох схем, тому нащадок може успадкувати вдалі фрагменти обох батьків.
\item \textit{Як мутація покращує глобальний пошук?}\\
Вона підтримує різноманітність і дозволяє перейти до схем, яких неможливо отримати лише кросовером наявних особин.
\item \textit{Наведіть приклади інженерної оптимізації.}\\
Розміщення компонентів, налаштування регуляторів, проєктування конструкцій, планування виробництва й оптимізація мереж.
\item \textit{Як ГА застосовують у маршрутизації та логістиці?}\\
Хромосома кодує маршрут або розподіл ресурсів, а fitness враховує відстань, час, вартість і обмеження.
\item \textit{Яке місце ГА серед методів глобальної оптимізації?}\\
Це популяційний еволюційний метод поряд із роєм частинок; на відміну від імітаційного відпалу, він одночасно підтримує багато кандидатів.
\item \textit{Як параметри ГА впливають на якість?}\\
Розмір популяції визначає покриття простору, кросовер поширює корисні комбінації, а мутація підтримує різноманітність.
\item \textit{Яка перевага ГА в багатокритеріальних задачах?}\\
Вони можуть працювати зі зваженою функцією або будувати набір Парето-оптимальних компромісів.
\item \textit{У яких реальних сферах використовують генетичну оптимізацію?}\\
У транспорті, телекомунікаціях, енергетиці, машинобудуванні, фінансах, біоінформатиці та машинному навчанні.
\end{enumerate}
\end{document}
